\documentclass[10pt]{article}

\usepackage[margin=0.9in]{geometry}
\usepackage{amsmath,amssymb,amsthm,mathtools}
\allowdisplaybreaks

\newtheorem{theorem}{Theorem}
\newtheorem{lemma}{Lemma}[section]

\newcommand{\defc}{\operatorname{defc}}
\newcommand{\NRCM}{\operatorname{NRCM}}
\newcommand{\clamp}{\operatorname{clamp}}
\newcommand{\ind}{\mathbf 1}

\title{Deficit Preservation for the Naive Rational Cyclic Map}
\author{}
\date{}

\begin{document}
\maketitle

\section{Definitions and statement}

Fix coprime positive integers \(r,s\), put \(\rho=r/s\), and index the
columns by \(0,1,\ldots,s-1\).  Write
\[
  H_i=\left\lfloor\frac{ri}{s}\right\rfloor,
  \qquad
  L_i\equiv ri\pmod s,
  \qquad 0\le L_i<s,
  \qquad
  \ell_i=\frac{L_i}{s}.
\]
The labels \(L_0,\ldots,L_{s-1}\) are distinct, and \(\ind_A\) denotes the
indicator of a condition \(A\).

A \emph{rational position-coordinate path} is an integer sequence
\(Q=(Q_0,\ldots,Q_{s-1})\) such that
\[
  Q_0=0,
  \qquad
  0\le Q_i\le H_i.
\]
Set
\[
  P_i=H_i-Q_i,
  \qquad
  q_i=Q_i+\ell_i.
\]
The path is \emph{valid} if
\[
  P_0\le P_1\le\cdots\le P_{s-1}.
\]
Since \(H_i+\ell_i=\rho i\), we shall repeatedly use
\begin{equation}
  q_i=\rho i-P_i.
  \label{eq:q-rho-P}
\end{equation}

For \(1\le i\le s-2\), define
\[
  C_i(Q)=[P_{i-1}-P_i,\,P_{i+1}-P_i].
\]
If \(I=[a,b]\), let \(\clamp_I(x)=a\) for \(x<a\),
\(\clamp_I(x)=x\) for \(a\le x\le b\), and
\(\clamp_I(x)=b\) for \(x>b\).  For \(1\le i<j<s\), put
\[
  \delta_{ij}(Q)
  =\left\lfloor
      \left|\clamp_{C_i(Q)}(q_i-q_j)\right|
    \right\rfloor,
\]
and define
\[
  \defc(Q)=
  \sum_{i=1}^{s-2}\sum_{j=i+1}^{s-1}\delta_{ij}(Q).
\]
Thus \(\defc(Q)=0\) when \(s\le2\).

For \(1\le k<s\), let \(I_k=\{k,k+1,\ldots,s-1\}\).  List the columns of
\(I_k\) in increasing label order as \((i_1,\ldots,i_{s-k})\), and define
\(T_k(Q)\) by fixing every coordinate outside \(I_k\) and setting
\[
  T_k(Q)_{i_{a+1}}=Q_{i_a}\quad(1\le a<s-k),
  \qquad
  T_k(Q)_{i_1}=Q_{i_{s-k}}+1.
\]
Thus the suffix values move once around the cyclic label order, with \(1\)
added on the wrap.

If some \(T_k(Q)\) satisfies all capacity inequalities
\(0\le T_k(Q)_i\le H_i\), choose the least such \(k\).  If this first
capacity-valid candidate is also a valid path, set
\[
  \NRCM(Q)=T_k(Q).
\]
Otherwise \(\NRCM(Q)\) is undefined.

\begin{theorem}[NRCM deficit preservation]
\label{thm:nrcm}
If \(Q\) is valid and \(\NRCM(Q)\) is defined, then
\[
  \defc(\NRCM(Q))=\defc(Q).
\]
\end{theorem}

\section{Proof}

If \(s=1\), there is no candidate suffix, so the theorem is vacuous.  Assume
\(s>1\), write
\[
  r=ms+d,
  \qquad
  m=\lfloor\rho\rfloor,
  \qquad
  1\le d<s,
\]
and let
\[
  Q'=\NRCM(Q)=T_k(Q),
  \qquad
  K=I_k.
\]
Both \(Q\) and \(Q'\) are valid.

We first replace the deficit by a pair energy.  The suffix move raises the
coordinate sum by \(1\), so it will be enough to show that the pair energy
also rises by \(1\).  We do this by splitting the pairs into fixed-prefix,
suffix-suffix, and mixed pairs.

\subsection{The pair energy}

For \(x\in\mathbb R\), define
\[
  G(x)=\max\{0,\lfloor x\rfloor\},
  \qquad
  \psi(x)=G(x)-G(x-\rho).
\]
For \(x\ge0\),
\begin{equation}
  \psi(x)
  =\#\{v\in\mathbb Z_{>0}:x-\rho<v\le x\}
  =\sum_{v\ge1}\ind_{[v,v+\rho)}(x).
  \label{eq:psi-intervals}
\end{equation}
Set
\[
  E(Q)=\sum_{0\le i<j<s}\psi(|q_i-q_j|),
  \qquad
  S(Q)=\sum_{i=1}^{s-1}Q_i.
\]

\begin{lemma}[Energy identity]
\label{lem:energy}
For every valid path \(Q\),
\begin{equation}
  \defc(Q)=E(Q)-S(Q).
  \label{eq:energy-identity}
\end{equation}
\end{lemma}

\begin{proof}
Put \(w_i=P_{i+1}-P_i\).  Since
\(C_i(Q)=[-w_{i-1},w_i]\), the definition of the clamp gives
\[
\begin{aligned}
\delta_{ij}(Q)
={}&G(q_i-q_j)-G(q_i-q_j-w_i)\\
 &+G(q_j-q_i)-G(q_j-q_i-w_{i-1}).
\end{aligned}
\]
By \eqref{eq:q-rho-P},
\[
  q_i-w_i=q_{i+1}-\rho,
  \qquad
  q_i+w_{i-1}=q_{i-1}+\rho,
\]
so
\begin{equation}
\begin{aligned}
\delta_{ij}(Q)
={}&G(q_i-q_j)-G(q_{i+1}-q_j-\rho)\\
 &+G(q_j-q_i)-G(q_j-q_{i-1}-\rho).
\end{aligned}
\label{eq:delta-G}
\end{equation}
Because the labels are distinct, \(q_i\ne q_j\) for \(i\ne j\), and hence
\begin{equation}
\begin{aligned}
\psi(|q_i-q_j|)
={}&G(q_i-q_j)+G(q_j-q_i)\\
 &-G(q_i-q_j-\rho)-G(q_j-q_i-\rho).
\end{aligned}
\label{eq:pair-G}
\end{equation}

Sum \eqref{eq:delta-G} over \(1\le i<j<s\) and compare with the sum of
\eqref{eq:pair-G} over \(0\le i<j<s\).  After reindexing the two negative
terms in \eqref{eq:delta-G}, the only terms left in \(E(Q)-\defc(Q)\) are
\[
\begin{array}{c|c}
\text{surviving term}&\text{range}\\ \hline
G(-q_j)+G(q_j)-G(-q_j-\rho)&1\le j<s\\
-G(q_1-q_j-\rho)&2\le j<s\\
-G(q_{i+1}-q_i-\rho)&0\le i<s-1.
\end{array}
\]
(The diagonal terms introduced by the reindexing are \(G(-\rho)=0\).)
For \(j>0\), coprimality gives \(0<\ell_j<1\), so
\[
  G(q_j)=Q_j,
  \qquad
  G(-q_j)=G(-q_j-\rho)=0.
\]
Moreover,
\[
  q_1-q_j-\rho=P_j-P_1-\rho j<0\quad(j\ge2),
\]
because \(P_j\le H_j<\rho j\), and
\[
  q_{i+1}-q_i-\rho=-(P_{i+1}-P_i)\le0
\]
by validity.  Thus only \(\sum_{j=1}^{s-1}Q_j=S(Q)\) remains, proving
\eqref{eq:energy-identity}.
\end{proof}

The suffix move permutes its \(Q\)-values and adds \(1\) once; hence
\begin{equation}
  S(Q')=S(Q)+1.
  \label{eq:S-increase}
\end{equation}
By Lemma~\ref{lem:energy}, it remains to prove
\[
  E(Q')=E(Q)+1.
\]

\subsection{Pairs inside the moving suffix}

List the columns of \(K\) in increasing label order as
\[
  a_1,\ldots,a_n,
  \qquad
  \ell_{a_1}<\cdots<\ell_{a_n},
  \qquad n=s-k,
\]
and put
\[
  u_j=\ell_{a_j}\quad(1\le j\le n),
  \qquad
  a_{n+1}=a_1,
  \qquad
  u_{n+1}=u_1+1.
\]
The value whose source is \(a_j\) moves from
\begin{equation}
  x_j=Q_{a_j}+u_j
  \quad\text{to}\quad
  y_j=Q_{a_j}+u_{j+1}.
  \label{eq:lifted-motion}
\end{equation}
Thus \(y_j>x_j\), and \(\{y_j:1\le j\le n\}\) is the multiset
\(\{q'_t:t\in K\}\).  This convention also covers \(n=1\): the sole point
moves from \(q_{a_1}\) to \(q_{a_1}+1\).

\begin{lemma}[Suffix-pair invariance]
\label{lem:suffix-pairs}
The contribution to \(E\) from pairs contained in \(K\) is unchanged by the
suffix move.
\end{lemma}

\begin{proof}
There is nothing to prove when \(n=1\).  Fix sources
\(a=a_j\) and \(c=a_\ell\) with \(j<\ell\), and let
\(a^+=a_{j+1}\), \(c^+=a_{\ell+1}\), with the lifted convention above.
Put
\[
  D=Q_c-Q_a,
  \qquad
  p=s(u_\ell-u_j),
  \qquad
  p'=s(u_{\ell+1}-u_{j+1}).
\]
The old and new contributions of this source pair are
\[
  \psi\!\left(\left|D+\frac ps\right|\right),
  \qquad
  \psi\!\left(\left|D+\frac{p'}s\right|\right).
\]
A direct evaluation of \(G\) gives, for every integer \(D\) and
\(1\le p<s\),
\begin{equation}
\psi\!\left(\left|D+\frac ps\right|\right)
=
\begin{cases}
D, &0\le D\le m,\\[1mm]
m+\ind_{\{p<d\}}, &D\ge m+1,\\[1mm]
-D-1, &-m-1\le D\le-1,\\[1mm]
m+\ind_{\{p>s-d\}}, &D\le-m-2.
\end{cases}
\label{eq:psi-pair}
\end{equation}
Thus a change is possible only when one of the thresholds \(d\) or
\(s-d\) is crossed.

We record the elementary geometry of such a crossing.  Let
\[
  x=su_j,
  \quad x^+=su_{j+1},
  \quad z=su_\ell,
  \quad z^+=su_{\ell+1};
\]
then \(x<x^+\le z<z^+\), \(p=z-x\), and \(p'=z^+-x^+\).  A physical step
one column to the right adds \(d\) to a lifted label.

If \(p<d\le p'\), then
\[
  z<x+d<x^++d\le z^+.
\]
Were \(a+1\in K\), its lifted label \(x+d\) would lie strictly between
\(c\) and its label successor \(c^+\).  Hence \(a=s-1\).  Now
\(a^++1\in K\), and its lifted label \(x^++d\) lies in \((z,z^+]\); it
must therefore equal \(z^+\).  Thus \(p'=d\) and \(c^+=a^++1\).
If instead \(p'<d\le p\), then
\[
  x<z^+-d<x^+.
\]
Were \(c^+-1\in K\), its lifted label would lie strictly between \(a\) and
\(a^+\); hence \(c^+=k\).  Moreover,
\[
  x\le z-d<z^+-d<x^+.
\]
Since \(c-1\in K\), its lifted label cannot lie strictly between \(a\) and
\(a^+\), so \(z-d=x\).  Thus \(p=d\) and \(c=a+1\).  Applying these
two conclusions to the reversed pair, whose old and new gaps are
\(s-p\) and \(s-p'\), shows that a change in the truth of \(p>s-d\) forces
\begin{equation}
  p=s-d,\ a=c+1,
  \qquad\text{or}\qquad
  p'=s-d,\ a^+=c^++1.
  \label{eq:reverse-threshold}
\end{equation}
All \(+1\) relations here are physical adjacency, not cyclic indexing.

Suppose first that \(D\ge m+1\).  By \eqref{eq:psi-pair}, a change would
force either \(p=d\) and \(c=a+1\), or \(p'=d\) and
\(c^+=a^++1\).  Validity of \(Q\) in the first case, and of \(Q'\) in the
second, gives respectively
\[
  D+\frac ds=q_c-q_a\le\rho,
  \qquad\text{or}\qquad
  D+\frac ds=q'_{c^+}-q'_{a^+}\le\rho.
\]
Either inequality implies \(D\le m\), a contradiction.

If \(D\le-m-2\), a change would force one of the alternatives in
\eqref{eq:reverse-threshold}.  Validity of the corresponding path gives
\[
  -D-1+\frac ds\le\rho,
\]
which implies \(D\ge-m-1\), again a contradiction.  Hence every source
pair has the same energy before and after the move, and summing over the
source pairs proves the lemma.
\end{proof}

\subsection{Why every omitted column overflows}

For \(0\le t<k\), insert the label \(L_t\) into the cyclic label order on
\(K\), and let \(\alpha_t\) and \(\beta_t\) be its predecessor and successor
(they coincide when \(|K|=1\)).  The source \(\alpha_t\) crosses the lifted
label of \(t\) at
\[
  c_t=Q_{\alpha_t}+\ell_t+
      \ind_{\{L_t<L_{\alpha_t}\}}.
\]
Its capacity margin at column \(t\) is
\begin{equation}
  \mu_t=\rho t-c_t
  =H_t-Q_{\alpha_t}-\ind_{\{L_t<L_{\alpha_t}\}}.
  \label{eq:mu-def}
\end{equation}
Thus \(\mu_t<0\) says exactly that the incoming value would exceed \(H_t\).
In particular,
\begin{equation}
  \mu_0=-Q_{\alpha_0}-1<0.
  \label{eq:mu-zero}
\end{equation}

\begin{lemma}[Omitted-column overflow]
\label{lem:overflow}
For every \(0\le t<k\), one has \(\mu_t<0\).
\end{lemma}

\begin{proof}
We first derive two useful projections of the same margin.  The motion of
source \(\alpha_t\) ends at \(q'_{\beta_t}\), and its lifted distance from
label \(t\) to label \(\beta_t\) is the fractional part of
\(\rho(\beta_t-t)\).  Therefore
\begin{equation}
  \mu_t
  =P'_{\beta_t}-\left\lfloor\rho(\beta_t-t)\right\rfloor.
  \label{eq:projection-right}
\end{equation}
Also,
\[
  H_{\alpha_t}-H_t
  =\left\lfloor\rho(\alpha_t-t)\right\rfloor
   +\ind_{\{L_{\alpha_t}<L_t\}}.
\]
Substituting \(Q_{\alpha_t}=H_{\alpha_t}-P_{\alpha_t}\) into
\eqref{eq:mu-def}, and observing that the two strict label indicators sum
to \(1\), gives
\begin{equation}
  \mu_t
  =P_{\alpha_t}-\left\lceil\rho(\alpha_t-t)\right\rceil.
  \label{eq:projection-left}
\end{equation}
Here \(0<\alpha_t-t<s\), so \(\rho(\alpha_t-t)\notin\mathbb Z\).

Equation \eqref{eq:mu-zero} settles the case \(k=1\).  Suppose \(k>1\), and
put
\[
  b=k-1,
  \qquad
  n=s-k.
\]
We compare every omitted column with \(b\).  For an integer \(g\), let
\(R(g)\) be the least nonnegative residue of \(dg\) modulo \(s\).  The
residues used below are nonzero and distinct.  Let
\(f,p\in\{1,\ldots,n\}\) minimize and maximize \(R(t)\), respectively.
For \(i=b-D\), where \(0\le D\le b-1\), let \(f_D,p_D\) minimize and
maximize \(R(D+t)\) on \(1\le t\le n\).  Relative to column \(i\), the
suffix column \(b+t\) has label residue \(R(D+t)\); hence
\begin{equation}
  \beta_i=b+f_D,
  \qquad
  \alpha_i=b+p_D,
  \qquad
  \beta_b=b+f,
  \qquad
  \alpha_b=b+p.
  \label{eq:alpha-beta-window}
\end{equation}

We need only the following window facts:
\begin{equation}
\begin{aligned}
  D+f_D&\ge f, & D+p_D&\ge p,\\
  f_D>f&\Longrightarrow p_D<p.
\end{aligned}
\label{eq:window-facts}
\end{equation}
For the first line, if the new minimum occurred at an absolute position
\(D+f_D<f\), then both it and \(f\) would lie in the shifted window, contrary
to the minimality of \(f_D\); the argument for the maximum is identical.

For the implication in the second line, order \(1,\ldots,n\) cyclically by
increasing \(R\)-value.  Its wrap edge is \(p\to f\).  Addition of \(R(D)\)
rotates this cyclic order, so \(p_D\to f_D\) is one of its edges.  We claim
that every ordinary edge \(a\to c\) with \(a>p\) has \(c<f\).  First,
\(f+p>n\): otherwise \(R(f+p)\), obtained from \(R(f)+R(p)\) modulo \(s\),
would be larger than the maximum, smaller than the minimum, or zero.  Now
suppose \(a>p\) and \(c\ge f\).  Since the edge is ordinary, \(c>f\).  Put
\(\Delta=R(c)-R(a)>0\).  If \(c>a\), then
\(1\le c-a<n-p<f\), so minimality of \(R(f)\) gives
\(\Delta=R(c-a)>R(f)\).  If \(c<a\), the position \(f+a-c\) lies in
\([1,n]\), and its residue is congruent to \(R(f)-\Delta\); minimality again
forces \(\Delta>R(f)\).  Consequently
\[
  R(a)<R(c)-R(f)=R(c-f)<R(c),
\]
contradicting the consecutiveness of \(a\to c\).  The claim follows.  If
\(f_D>f\), the edge \(p_D\to f_D\) is therefore sourced at or left of \(p\);
it cannot be sourced at \(p\), whose successor is \(f\).  Hence \(p_D<p\),
proving \eqref{eq:window-facts}.

By \eqref{eq:alpha-beta-window} and \eqref{eq:window-facts},
\begin{equation}
  \beta_i-i\ge\beta_b-b,
  \qquad
  \alpha_i-i\ge\alpha_b-b,
  \qquad
  \beta_i>\beta_b\Longrightarrow\alpha_i<\alpha_b.
  \label{eq:suffix-window}
\end{equation}
If \(\beta_i\le\beta_b\), validity of \(Q'\),
\eqref{eq:projection-right}, and \eqref{eq:suffix-window} give
\[
  \mu_i
  \le P'_{\beta_b}-\left\lfloor\rho(\beta_b-b)\right\rfloor
  =\mu_b.
\]
If \(\beta_i>\beta_b\), then \(\alpha_i<\alpha_b\), so validity of \(Q\),
\eqref{eq:projection-left}, and \eqref{eq:suffix-window} give
\[
  \mu_i
  \le P_{\alpha_b}-\left\lceil\rho(\alpha_b-b)\right\rceil
  =\mu_b.
\]
Thus
\begin{equation}
  \mu_i\le\mu_b\qquad(1\le i\le b).
  \label{eq:margin-dominance}
\end{equation}

It remains to show \(\mu_b<0\).  If \(\mu_b\ge0\), insert column \(b\) into
the selected suffix cycle.  The incoming value from \(\alpha_b\) fits at
\(b\) by \eqref{eq:mu-def}.  The only other changed target is \(\beta_b\),
which receives
\[
  Q_b+\ind_{\{L_{\beta_b}<L_b\}}.
\]
This value also fits, because
\[
  H_{\beta_b}-H_b
  =\left\lfloor\rho(\beta_b-b)\right\rfloor
   +\ind_{\{L_{\beta_b}<L_b\}}
  \ge \ind_{\{L_{\beta_b}<L_b\}},
\]
and \(Q_b\le H_b\).  Every other target receives exactly what it received
under the capacity-valid \(K\)-move.  Hence \(T_b(Q)\) would be
capacity-valid, contradicting the minimality of \(k\).  Therefore
\(\mu_b<0\), and \eqref{eq:margin-dominance}, together with
\eqref{eq:mu-zero}, proves the lemma.
\end{proof}

\subsection{The mixed pairs}

Extend the preceding notation by setting
\[
  \mu_k=P'_k,
  \qquad
  c_k=q'_k=\rho k-\mu_k.
\]
Thus \(c_t=\rho t-\mu_t\) for every \(0\le t\le k\).

\begin{lemma}[One-turn mixed-energy count]
\label{lem:mixed}
The contribution to \(E\) from pairs with one endpoint in
\(\{0,\ldots,k-1\}\) and the other in \(K\) increases by exactly \(1\).
\end{lemma}

\begin{proof}
For \(0\le h<k\), write
\[
  F_h(Q)=\sum_{j\in K}\psi(|q_j-q_h|).
\]
By \eqref{eq:psi-intervals}, for fixed \(q_h\),
\begin{equation}
\begin{aligned}
\psi(|x-q_h|)=\sum_{v\ge1}\bigl(
 &\ind_{[q_h+v,\,q_h+v+\rho)}(x)\\
 &+\ind_{(q_h-v-\rho,\,q_h-v]}(x)
\bigr).
\end{aligned}
\label{eq:interval-decomposition}
\end{equation}
Move each suffix point from \(x_j\) to \(y_j\) as in
\eqref{eq:lifted-motion}.  For \([A,B)\), entrance and exit crossings satisfy
\(x_j<A\le y_j\) and \(x_j<B\le y_j\); for \((A,B]\), they satisfy
\(x_j\le A<y_j\) and \(x_j\le B<y_j\).  Entrances contribute \(+1\) and
exits \(-1\).  The label arcs from \(u_j\) to \(u_{j+1}\) partition one
full turn.  Thus an omitted label \(t<k\) is crossed exactly once, at
\(c_t\); at the terminal suffix label \(k\), the right-interval convention
selects the endpoint \(c_k=q'_k\).  This remains true for a singleton suffix.

Assume first that \(h\ge1\).  The four endpoint families in
\eqref{eq:interval-decomposition}, their signs, and the corresponding
positive-integer crossing conditions are
\[
\begin{array}{c|c|c|c}
\text{endpoint}&\text{sign}&\text{label}&\text{required }v\\ \hline
q_h-v-\rho&+&h-1&\mu_{h-1}-P_h\\
q_h+v&+&h&P_h-\mu_h\\
q_h-v&-&h&\mu_h-P_h\\
q_h+v+\rho&-&h+1&P_h-\mu_{h+1}.
\end{array}
\]
Indeed, each entry in the last column follows at once from
\(q_h=\rho h-P_h\) and \(c_t=\rho t-\mu_t\), and it contributes precisely
when that entry is a positive integer.

For \(h=0\), all left intervals lie on the negative half-line, while every
suffix point stays positive.  Exactly one right entrance is crossed, at the
positive integer
\[
  c_0=-\mu_0=Q_{\alpha_0}+1,
\]
so it contributes \(+1\).  A right exit is crossed exactly when
\(c_1=\rho+v\), equivalently \(v=-\mu_1>0\).

If \(k=1\), validity of \(Q'\) gives
\(\mu_1=P'_1\ge P'_0=0\).  Hence no right exit is crossed at \(h=0\), and
the mixed energy rises by \(1\).

Now suppose \(k>1\).  Lemma~\ref{lem:overflow} gives
\[
  \mu_0,\ldots,\mu_{k-1}<0,
\]
whereas \(P_h\ge0\).  At \(h=0\), both the right entrance and the right exit
are crossed, so the net change is \(0\).  For
\(1\le h\le k-2\), the table shows that exactly the second and fourth
crossings occur, again giving net change \(0\).  At \(h=k-1\), the second
crossing still occurs, but the fourth does not, because validity of \(Q'\)
gives
\[
  \mu_k=P'_k\ge P'_{k-1}=P_{k-1}.
\]
Thus \(F_{k-1}(Q')-F_{k-1}(Q)=1\).  Summing over the fixed prefix proves the
lemma.
\end{proof}

\begin{proof}[Proof of Theorem~\ref{thm:nrcm}]
Pairs entirely in the fixed prefix do not move.  By
Lemma~\ref{lem:suffix-pairs}, the suffix-suffix energy is unchanged, while
Lemma~\ref{lem:mixed} shows that the mixed energy rises by \(1\).  Hence
\[
  E(Q')-E(Q)=1.
\]
Using Lemma~\ref{lem:energy} and \eqref{eq:S-increase}, we conclude that
\[
\begin{aligned}
\defc(Q')-\defc(Q)
&=\bigl(E(Q')-E(Q)\bigr)-\bigl(S(Q')-S(Q)\bigr)\\
&=1-1=0.
\end{aligned}
\]
Therefore \(\defc(\NRCM(Q))=\defc(Q)\).
\end{proof}

\end{document}
